\Problem{Consulta de Percentiles}

\Statement

Disponemos de los datos de un ranking de jugadores, cada uno con su nombre (un \verb|string|) y sus puntos (un \verb|int|). La puntuación de todos los jugadores es única (no hay empates). Si los ordenamos por puntuación de mayor a menor, el \textbf{percentil} de un jugador es el porcentaje de jugadores que tiene por debajo (incluyéndose a sí mismo) en el ranking.

\medskip

Por ejemplo, en el siguiente ranking se muestran el nombre, los puntos y el percentil de un grupo de 5 jugadores, ordenados por puntos descendentemente:

\medskip

\begin{tabular}{lcc}
    \verb|ramon| & 15 & 100\% \\
    \verb|amy| & 13 & 80\% \\
    \verb|john| & 10 & 60\% \\
    \verb|pete| & 8 & 40\% \\
    \verb|zach| & 4 & 20\% \\
\end{tabular}

\medskip

Haced un programa que lea las puntuaciones de un ranking de jugadores y después conteste una serie de consultas, en las que se pide un nombre y hay que contestar con el percentil del jugador pedido. El programa se debe diseñar teniendo en cuenta que el número de consultas es, típicamente, \emph{dos órdenes de magnitud mayor} que el número de jugadores del ranking.

\Observation

En este problema el centro de interés es la eficiencia. Se trata de aplicar los mejores algoritmos conocidos y encontrar formas de hacer el mínimo número de operaciones para poder producir la salida.

\medskip

El uso del algoritmo \verb|sort| de \verb|vector| está permitido en este problema. Como recordatorio, para usar \verb|sort|, haría falta:
\begin{itemize}
\item Añadir \verb|#include <algorithm>|.
\item Definir una función de comparación \verb|fcomp| que reciba dos elementos constantes \verb|a| y \verb|b| y retorne \verb|true| si \verb|a| va antes que \verb|b| en la ordenación final.
\item Llamar a \verb|sort(v.begin(), v.end(), fcomp)|, suponiendo que \verb|v| es el vector que se quiere ordenar.
\end{itemize}

\Input

La entrada tiene dos partes. Primero hay una secuencia de nombres de jugadores, acabada en "\verb|---|", donde cada jugador va acompañado de su puntuación. Las puntuaciones de todos los jugadores son diferentes, es decir, no hay jugadores empatados a puntos. El orden de los jugadores es arbitrario. Después sigue una secuencia sin centinela de nombres de jugadores de los cuales se quiere consultar el percentil.

\Output

La salida es simplemente el percentil de cada jugador consultado, uno por línea. Los percentiles hay que mostrarlos con 2 decimales. Para mostrar dos decimales con todos los números reales de un programa se suele poner, al principio de la función \verb|main|:

\begin{verbatim}
cout.precision(2);
cout.setf(ios::fixed);
\end{verbatim}

\SampleTwoCol
